\documentstyle[%


righttag%


,11pt%


,amssymb%


,russian%               remove in Eng.


,izvru%                 Set izven% in Eng.


]{amsart}





%


\newcommand{\la}{\langle}


\newcommand{\ra}{\rangle}


\newcommand{\gr}{\text{\,гр\,}} %% repl for   \mathrm{gr}   in eng.


\newcommand{\diag}{\operatorname{diag}}


\newcommand{\tts}[1]{} 


\renewcommand{\baselinestretch}{0.99}








%\hoffset=-15mm%                  For Eng.


\setcounter{page}{51}


\god{2003}


\nomer{10~(497)} %                        Remove in Eng.


%\nomer{Vol.\,47, No.\,10}%               Set in Eng.


%\str{\pageref{begin}--\pageref{end}}%  Set in Eng.


\udk{512.552}





\author{\it А.М.\,Попова }


% NB:   ^^ remove in Eng.


\title{Мультипликативная структура конечно-порожденных


матричных колец}


\thanks{Работа выполнена при финансовой поддержке Российского фонда


фундаментальных исследований (код проекта \No\,99-01-00-571).}





\date{}


%\pagestyle{myheadings}%         Set in Eng.


\pagestyle{plain} %              Remove in Eng.


\begin{document}


%\thispagestyle{plain}%          Set in Eng.


\maketitle


\label{begin}











\section{Постановка задачи}





Описание групп единиц различных ассоциативных колец с единицей и нахождение


условий на группу, при которых она является мультипликативной группой


подходящего кольца, --- эти две проблемы были сформулированы Л.\,Фуксом в


(\cite{Fuch1}, гл.XIII, \S\,77, с.\,299; \cite{Fuch2},


гл.\,XVIII, \S\,129, с.\,380). В данной работе эти задачи


решаются для произвольного конечно-порожденного матричного кольца над полем


рациональных чисел.





Полученные результаты позволяют решать задачу описания групп 


единиц целочисленных групповых колец конечных групп и групп, представимых


матрицами (\cite{Good}--\cite{Al}).





Пусть $O=rg(x_{1},\dotsc,x_{r})\subset Q_{n}$ --- такое кольцо. Прежде всего,


заметим, что $O$ приводится над $Q$ к клеточно-треугольному виду, в котором


диагональные клетки либо абсолютно неприводимы,


либо неприводимы над $Q$. Поэтому естественным подходом к решению поставленных


задач является вначале решение их для каждой клетки в отдельности, а потом


``склейка'' полученных результатов для общего случая.





В работе \cite{Pop1} доказывается эффективность приведения $O$ к


клеточно-треугольному виду, т.\,е. эффективность нахождения по порождающим


$x_1,\dotsc,x_r$ кольца $O$ порождающих каждой клетки, а один из результатов


работы \cite{Pop2} позволяет эффективно отвечать на вопрос, содержит


ли клетка единичную матрицу.





Пусть $O_i$ --- такая клетка. Обозначим


через $QO_i$ алгебру кольца $O_i$ над полем рациональных чисел.


Могут возникнуть три случая:





1) $QO_i = Q_{n_i}$;





2) $QO_i \cong F_k$, $F$ --- поле;





3) $QO_i \cong T_k$, $T$ --- тело.





\noindent{}Первый случай относится к абсолютно 


неприводимым кольцам, второй и третий ---


к кольцам, неприводимым над $Q$. Для первых двух случаев задача описания группы


единиц $U(O_i)$ кольца $O_i$ на языке порождающих элементов была решена в


\cite{Pop3} и \cite{Pop4}. Третий случай рассмотрен в \cite{Pop5}.


Таким образом, в данной работе поставленные задачи решаются уже для 


произвольного конечно-порожденного матричного кольца.





В силу вышеизложенного будем полагать, что $O$ имеет клеточно-треугольный вид


и содержит единичную матрицу


\begin{equation}


O=


\begin{pmatrix}


\framebox{$O_1$} & & * \\ 


 & \ddots & \\ 


 0 & & \framebox{$O_q$}


\end{pmatrix}


.


\end{equation}





\section{Строение абсолютно неприводимых и неприводимых над $Q$ колец}








В силу того, что в случае неприводимости кольца $O$ алгебра $QO$ 


изоморфна полной матричной алгебре над полем или телом (\cite{Sup},


с.\,142), справедлива 








\begin{thm}


Если конечно-порожденное кольцо $O \subset Q_n$ абсолютно


неприводимо или неприводимо над $Q$, то эффективно находятся такие элементы


$g_1,\dotsc,g_t \in O$ и такое множество 


простых чисел $\pi = \la p_1,\dotsc,p_s\ra$,


что


$$


O = \{ g_1,\dotsc,g_t \}_{Z_{\pi}},


$$


где $Z_{\pi} = Z[\frac{1}{p_1 \dotsb p_s}]$.


\end{thm}








{\bf Доказательство} теоремы можно найти в \cite{Pop2} 


для абсолютно неприводимого


случая, аналогично она может быть доказана для кольца,


неприводимого над $Q$. Для неприводимого кольца будем


считать, что $QO=H_k$, где


$H\subset Q_l$ --- тело размерности $l=\frac{n}{k}$ 


(\cite{Sup}, с.\,142).


Обозначим $H_{\pi}=H\cap (Z_{\pi})_l$. 


Из теоремы 1 легко следует, что найдется


ненулевое натуральное число $m$ такое, что $(m,p_1\cdots p_s)=1$ и


справедливы следующие включения:





$m(Z_{\pi})_n\subset O$ для абсолютно неприводимого случая,


$m(H_{\pi})_k \subset O$ для неприводимого случая.





В кольцах $Z_{\pi}$ и $H_{\pi}$ породим соответственно идеалы $I=(m)$


и $I'=(me_l)$ ($e_l$ --- единичная матрица из $Q_l$). Определим гомоморфизмы


\medskip





\begin{center}


\begin{tabular}{ll}


абсолютно неприводимый случай & неприводимый случай \\[2mm]


$\gamma_m : Z_{\pi} \longrightarrow Z_{\pi}/I=\overline{Z_{\pi}}$, & 


$\gamma '_m: H_{\pi} \longrightarrow H_{\pi}/I'= \overline{H_{\pi}}$,\\[2mm]


$\psi_m: Z_{\pi^n}\longrightarrow \overline{Z_{\pi}}_n$, & 


$\psi'_m: H_{\pi ^k} \longrightarrow \overline{H_{\pi}}_k$, \\[2mm]


$\varphi_m : GL_n(Z_{\pi})\longrightarrow GL_n(\overline{Z_{\pi}})$, & 


$\varphi '_m : GL_k(H_{\pi})\longrightarrow GL_k(\overline{H_{\pi}})$, \\[2mm]


$U_0(O)=U(O)\cap \ker\varphi_m$; & $U_0(O)=U(O)\cap \ker\varphi '_m$.


\end{tabular}


\end{center}








\section{Группы единиц абсолютно неприводимых и неприводимых над $Q$ колец}





Введем некоторые обозначения. Для абсолютно неприводимого случая


\begin{gather}


\gamma^*_m : U(Z_{\pi})\longrightarrow U(\overline{Z_{\pi}}),\ \ 


\ker\gamma^*_m = \gr(\rho_1,\dotsc,\rho_q);\ \


t_{ij}(\lambda ) = e + \lambda e_{ij},\ \


d(\rho ) = \diag(1,\dotsc,1,\rho ), \notag \\


%


U_0(m,\pi) = \gr(t_{ij}(\tfrac{m}{p_1\dotsb p_s}),


\ d(\rho_{\beta}),\ i,j=1,\dotsc,n,


\ \beta =1,\dotsc,q).


\end{gather}





Для неприводимого случая $QO=H_k$, где $H$ --- поле. Поскольку


$O\subset H_k$, то элементы $O$ имеют


клеточное строение, при котором матрица размера $n\times n$ разбита на


$k^2$ клеток размера $l\times l$.


В связи с этим пусть $e_{ij}$ --- матрица из $Q_n$, в которой в клетке


с номером $ij$ ($i,j=1,\dotsc,k$)


стоит единичная матрица из $Q_l$, $E_k=e_{11}+


e_{22}+\dotsb +e_{kk}$. Для $\mu \in H$, $g\in Q_n$ обозначим


\begin{gather}


\mu E_k = \diag(\mu,\dotsc,\mu ), \ \


\mu g=\mu E_kg,\ \ 


t_{ij}(\mu )=E_k+\mu e_{ij},\ \ 


d(\mu )=\diag(e_l,\dotsc,e_l,\mu );\notag \\


{\gamma '}_m^* : U(H_{\pi})\longrightarrow U(\overline{H_{\pi}}),\ \ 


\ker{\gamma '}_m^*=\gr(\tau_1,\dotsc,\tau_v); \notag \\


%


U_0(m,\pi,H_{\pi})=\gr(t_{ij}(\tfrac{mh}{p_1\dotsb p_s}),


\ d(\tau_{\beta}),


\ i,j=1,\dotsc,k,\ \beta =1,\dotsc,v).


\end{gather}


Пусть теперь $O$ неприводимо над $Q$ и $QO=H_k$, где $H$ --- тело. 


Известно (\cite{Jac}, c.\,266), что в этом случае


$[H:Q]=d^2$ и в $H$ существуют два подполя $F=\{1,f,\dotsc,f^{d-1}\}_Q$ 


и $T=\{1,t,\dotsc,t^{d-1}\}_Q$ такие, что





1) $H=\{f^it^j|i,j=0,\dotsc,d-1\}_Q$,





2) $\exists \xi \in H[f={\xi}^{-1}t\xi ]$.





\noindent{}Отсюда следует, что 


для любого $h\in H$ справедливо представление $h=ab$,


где $a\in F$, $b\in T$. Тогда, как показано в \cite{Pop5},


$H_{\pi}=F_{\pi}T_{\pi}$.


Поскольку $F$ и $T$ --- алгебраические расширения $Q$, то для максимальных


порядков в них существует теория дивизоров (\cite{Bor}, c.\,216),


с помощью которой можно найти порождающие групп $U(F_{\pi})$ и $U(T_{\pi})$


и тем самым порождающие группы $U(H_{\pi})=U(F_{\pi})U(T_{\pi})$.


В \cite{Pop5} доказывается, что порождающие


группы $\ker \gamma'{}^*_m$ и в этом


случае находятся эффективно. Пусть


$\ker \gamma'{}^*_m=\gr(\eta_1,\dotsc,\eta_w)$.


Обозначим


\begin{equation}


U_0(m,\pi,F_{\pi},T_{\pi})=\gr(t_{ij}(\tfrac{mf}{p_1\dotsb p_s}),


t_{ij}(\tfrac{mt}{p_1\dotsb p_s}),\ d(\eta_{\beta}),


\ i,j=1,\dotsc,k,\ \beta =1,\dotsc,w).


\end{equation}


В работах \cite{Pop3} и \cite{Pop6} доказаны следующие теоремы.








\begin{thm}


Абсолютно неприводимая группа $G< GL_n(Q)$ $(n>2)$ является группой


единиц конечно-порожденного абсолютно неприводимого матричного кольца


над полем рациональных чисел $Q$ тогда и только тогда, когда





$1)$ $G$ есть конечное расширение $U_0(m,\pi )$ для некоторых $m$ и $\pi;$





$2)$ если $Z[G]$ --- целочисленная линейная оболочка $G$, то $\overline{U(Z[G])}=


\overline{G}$, где черта означает образ при гомоморфизме $\varphi_m$.


\end{thm}








\begin{thm}


Неприводимая над $Q$ группа $G< GL_n(Q)$ $(n>2)$ является группой


единиц конечно-порожденного неприводимого над $Q$ матричного кольца 


тогда и только тогда, когда выполняются следующие условия\/${:}$





$1)$ существует $t\in GL_n(Q)$ такая, что $G^t\subset H_{\pi ^k}$ и $G^t$ есть


конечное расширение $U_0(m,\pi,H_{\pi})$ для некоторых $m$, $\pi$, $H_{\pi}$


или $U_0(m,\pi,F_{\pi},T_{\pi})$ для некоторых


$m$, $\pi$, $F_{\pi}$, $T_{\pi};$





$2)$ если $Z[G]$ --- целочисленная линейная оболочка $G$, то


$\overline{U(Z[G])}=\overline{G}$, где черта


означает образ при гомоморфизме $\varphi '_m$.


\end{thm}








Кроме того, в обоих случаях справедлива








\begin{thm}


Если $O=rg(x_1,\dotsc,x_r)\subset Q_n $ --- конечно-порожденное


абсолютно неприводимое или неприводимое над $Q$ кольцо, то задача нахождения


порождающих группы единиц $U(O)$


кольца $O$ алгоритмически разрешима.


\end{thm}








\section{Общий случай}








Если $O=rg(x_1,\dotsc,x_r)\subset Q_n $ вполне приводимо над $Q$ с


двумя неприводимыми частями $O_1$ и $O_2$ так, что $O=\diag (O_1,O_2)$, то


так же, как в \cite{Pop7}, можно показать, что либо существует такое


ненулевое $m\in N$, что $\diag (mO_1,0), \diag(0,mO_2)\subset O$, т.\,е.


$O_1$ и $O_2$ ``расклеиваются'', либо $O_1$ и $O_2$ не


``расклеиваются'', и тогда они изоморфны.





Так как в виде (1) кольца $O_i(i=1,\dotsc,q)$ определяются с точностью до


изоморфизма и порядка следования, то будем считать, что в последовательности


$O_1,\dotsc,O_q$ не ``расклеивающиеся'' кольца следуют


друг за другом и образуют блоки $R_1,\dotsc,R_{s+1}$, т.\,е.


\begin{equation} 


O=


\begin{pmatrix}


\framebox{$R_1$} & & * \\ 


& \ddots & \\ 


0 & & \framebox{$R_{s+1}$}


\end{pmatrix}


.


\end{equation}


Нетрудно показать, что $O=O'\oplus I$, где $O'\cong \diag(R_1,\dotsc,R_{s+1})$,


$I$ --- ядро гомоморфизма


\begin{equation}


\chi :


\begin{pmatrix}


\framebox{$R_1$} & & * \\ 


& \ddots & \\ 


0 & & \framebox{$R_{s+1}$}


\end{pmatrix}


\longrightarrow


\begin{pmatrix}


\framebox{$R_1$} & & 0 \\ 


& \ddots & \\ 


0 & & \framebox{$R_{s+1}$}


\end{pmatrix}


.


\end{equation}





Поскольку $U(O)=U(O')(e\oplus I)$, то естественно вначале описать группы


типа $U(R_i)$, затем типа $U(R)$, где $R=\diag(R_1,\dotsc,R_{s+1})$, и,


наконец, $U(O)$. Все это обосновано в \cite{Pop7} 


и дословно переносится на изученный здесь


случай. Поэтому приведем только формулировки необходимых лемм.








\begin{lem}


Пусть $ O=rg(x_1,\dotsc,x_r)\subset Q_n $ --- приводимое над $Q$ кольцо,


при этом неприводимые части либо аболютно неприводимы, либо неприводимы над 


$Q$ и в обоих случаях не


``расклеиваются''. Тогда эффективно находятся аддитивный базис


кольца $O$ и


конечная система порождающих группы единиц $U(O)$.


\end{lem}








Заметим, что если кольцо $O$ удовлетворяет условиям леммы 1, то


$U_0(O)=U(O)\cap \ker\varphi_m=(U(O')\cap \ker\varphi_m)(e\oplus mI)$


(аналогично для $\varphi '_m$),


поэтому группа


\begin{equation}


U_0(O)=\gr\left(


\begin{pmatrix}


\framebox{$u_{01}^1$} & & * \\ 


& \ddots & \\ 


0 & & \framebox{$u_{01}^t$}


\end{pmatrix}


,\dotsc,


\begin{pmatrix}


\framebox{$u_{0s}^1$} & & * \\ 


& \ddots & \\ 


0 & & \framebox{$u^t_{0s}$}


\end{pmatrix}


,e+ma_1,\dotsc,e+ma_q\right),


\end{equation}


где


$$


\gr\left(


\begin{pmatrix}


\framebox{$u_{01}^1$} & & * \\ 


& \ddots & \\ 


0 & & \framebox{$u_{01}^t$}


\end{pmatrix}


,\dotsc,


\begin{pmatrix}


\framebox{$u_{0s}^1$} & & * \\ 


& \ddots & \\ 


0 & & \framebox{$u^t_{0s}$}


\end{pmatrix}


\right)


= U_{01}


$$


изоморфна группе $U_{01}^1=\gr(u_{01}^1,\dotsc,u_{0s}^1)$, которая имеет вид


(2), (3) или (4), a $I=\{a_1,\dotsc,a_q\}_{Z_{\pi}}$.








\begin{lem}


Для любого вполне приводимого конечно-порожденного кольца над


полем рациональных чисел эффективно находится аддитивный базис.


\end{lem}








Вернемся к общему случаю. Пусть теперь $O$ имеет вид (5). Используя 


аддитивный базис кольца $R=\diag(R_1,\dotsc,R_{s+1})$, получим


$$


O=O'\oplus I,


$$


где $O'\cong R$, $I$ --- ядро гомоморфизма $O\longrightarrow O'$. 


Как следует


из замечания к лемме~1, для каждого $R_i$ эффективно находятся порождающие


групп $U_0(R_i)=U(R_i)\cap \ker\varphi_{im}$, поэтому легко,


применяя рассуждения такие же, как в \cite{Pop3} и \cite{Pop4}, 


эффективно найти


порождающие группы $U(R)$ и, следовательно, порождающие группы $U(O')$.


Поскольку, как и выше, $U(O)=U(O')(e\oplus I)$, то для нахождения конечной


системы порождающих $U(O)$ нужно любую матрицу из $e\oplus I$ представлять


словом от какого-то конечного набора матриц из $e\oplus I$ и матриц из $U(O')$.


Назовем этот конечный набор системой элементов,


задающих подгруппу $e\oplus I$, и обозначим его через $B$. 


Алгоритм построения множества


$B=\{ h_{11},\dotsc,h_{1l_1},\dotsc,h_{s1},\dotsc,


h_{sl_s}\} $ описан в \cite{Pop7}.








\begin{lem}


Пусть $h\in I$, тогда $e+h=w((e+h_{ij}),u_k)$, где $h_{ij}\in B$, 


$u_k\in U(O')$.


\end{lem}





\section{Основные результаты}








\begin{thm}


Пусть кольцо $O=rg(x_1,\dotsc,x_r)\subset Q_n$ $(n>2)$ приводимо над


полем рациональных чисел, при


этом неприводимые части либо абсолютно неприводимы, либо неприводимы 


над~$Q$. Тогда задача нахождения


порождающих группы единиц $U(O)$ алгоритмически разрешима.


\end{thm}








Поскольку $U(O)=U(O')(e\oplus I)$, то утверждение теоремы следует из


замеченного выше факта об эффективности нахождения порождающих $U(O')$ и


леммы 3. Таким образом, если $U(O')=\gr(u_1,\dotsc,u_t)$,


то 


$$


U(O)=\gr(u_1,\dotsc,u_t,e+h_{11},\dotsc,e+h_{1l_1},


\dotsc,e+h_{s1},\dotsc,e+h_{sl_s}).


$$





В следующей теореме сформулированы необходимые и достаточные условия для того, чтобы 


группа $G$ из $GL_n(Q)$ являлась группой единиц конечно-порожденного


матричного кольца над полем рациональных чисел.








\begin{thm}


Группа $G<GL_n(Q)$ $(n>2)$ является группой единиц конечно-порожденного 


матричного кольца над полем рациональных чисел $Q$ тогда и


только тогда, когда она


удовлетворяет следующим условиям\/${:}$





%% NB: repl. гр for \gr or \mathrm{gr} vvvvvvvv





$1)$ существует такая матрица $t\in GL_n(Q)$, что


$G^t=гр\,(u_1,\dotsc,u_t,e+h_1,\dotsc,e+h_q)$,


где группа $U_1=гр\,(u_1,\dotsc,u_t)$


изоморфна конечному расширению группы 


$$


U_0=гр\,\left(


\begin{pmatrix}


\framebox{$e_1$} & & & & * \\ 


& \ddots & & & \\


& & \framebox{$u_{ij}$} & & \\


& & & \ddots & \\ 


0 & & & & \framebox{$e_k$}


\end{pmatrix},\ \ i=1,\dotsc,k,\ \ j=1,\dotsc,t_i\right),


$$


для которой группы $U_{0i}=гр\,(u_{i1},\dotsc,u_{it})$ имеют вид


$(2)$--$(4)$ или $(7);$





$2)$ если $\chi$ --- гомоморфизм кольца $Z[G^t]$, определенный аналогично


$(6)$, $I=\ker\chi $, то система $\la h_1,\dotsc,h_q\ra$


аддитивных порождающих $I$ обладает свойством


$$


\forall h\in I[e+h=w((e+h_i),u_j),\ u_j\in U_1,\ j=1,\dotsc,q];


$$





$3)$ если $Z[G]$ --- целочисленная 


линейная оболочка $G$, то $\overline{U(Z[G])}


=\overline G$, где черта означает образ при гомоморфизме $\varphi_m$ 


$(\varphi '_m)$.


\end{thm}








\begin{thebibliography}{99}





\bibitem{Fuch1}


Fuchs L. {\em Abelian Groups}. -- Budapest: Acad.


Sci., 1958. -- 367p.





\bibitem{Fuch2}


Фукс Л. {\em Бесконечные абелевы группы}. Т.\,2. -- М.: Мир, 1977.


-- 416 с.





\bibitem{Good}


Goodaire E.G., Jespers E., Parmenter M.M. {\em Determining


units in some integral group rings} // Canad. Math. Bull.


-- 1990. -- V.\,33. -- \No\,2.


-- P.\,242--246.





\bibitem{Ritt}


Ritter J., Sehgal S.K. {\em Construction of units in integral


group rings of finite nilpotent groups} // Trans. Amer. Math.


Soc. -- 1991. --


V.\,324. -- \No\,2. -- P.\,603--621.





\bibitem{Al}


Алеев Р.Ж. {\em Единицы полей характеров и центральные единицы


целочисленных групповых колец конечных групп} // Матем. тр. -- Новосибирск,


2000. -- Т.\,3. -- \No\,1. --


С.\,3--37.





\bibitem{Pop1}


Попова А.М. {\em Некоторые алгоритмические проблемы для матричных


колец}. -- М., 1979. -- 9~с. -- Деп. в ВИНИТИ 27.07.79, \No\,2852-79.





\bibitem{Pop2}


Попова А.М. {\em Нахождение определяющих соотношений для


конечно-порожденного модуля над конечно-порожденным матричным кольцом}


-- М., 1984. -- 15 с. -- Деп. в ВИНИТИ 24.08.84, \No\,6010-84.





\bibitem{Pop3}


Попова А.М. {\em Группы единиц абсолютно неприводимых матричных


колец} // Актуальн.\linebreak пробл. современ. матем. -- Новосибирск, 1997. --


Т.\,3. -- С.\,155--161.





\bibitem{Pop4}


Попова А.М. {\em Мультипликативные группы конечно-порожденных


неприводимых матричных колец} // Фундаментальн. и прикл. матем.


-- 1999. -- Т.\,5. -- \No\,3. -- С.\,743--749.





\bibitem{Pop5}


Попова А.М. {\em Описание групп единиц неприводимых


матричных колец} // Алгебра и теория моделей. -- Новосибирск,


1997. -- С.\,152--159.


 


\bibitem{Sup}


Супруненко Д.А. {\em Группы матриц}. -- М.: Наука, 1972. -- 351~с.





\bibitem{Jac} 


Джекобсон Н. {\em Строение колец}. -- М.: Ин. лит., 


1961. -- 392 с.





\bibitem{Bor} 


Боревич З.И., Шафаревич И.Р. {\em Теория чисел}. -- М.: Наука, 1985.


-- 503 с.





\bibitem{Pop6}


Попова А.М. {\em Конечно-порожденные кoльца и их мультипликативная 


структура} // Тр. конф. по математике, посвящ. памяти А.Д.\,Тайманова.


-- Алматы, 1998. -- С.\,114--117.





\bibitem{Pop7}


Попова А.М. {\em Группы обратимых элементов матричных колец} // Сиб.


матем. журн. -- 1999. -- Т.\,40. -- \No\,5. -- С.\,1127--1136.





\end{thebibliography}





\vskip 0.5cm





\post{\it Новосибирский государственный}{\it Поступила}


\post{\it технический университет}{02.07.2001}





\label{end}


\end{document}








